\subsection{Pseudo-isomorphisms}

We preserve the notation and hypotheses of nos. 2 and 3.

PROPOSITION 9. Let M be a finitely generated A-module. The following conditions are equivalent:

\begin{enumerate}
\def\labelenumi{(\alph{enumi})}
\item
  \(\mathbf{M}_{\mathfrak{p}} = 0\) for every prime ideal \(\mathfrak{p}\) of height \(\leqslant 1\) .
\item
  The annihilator \(\mathfrak{a}\) of \(\mathbf{M}\) is an ideal \(\neq (0)\) and \(\mathbf{A}:\mathfrak{a} = \mathbf{A}\) ( \(\mathbf{A}:\mathfrak{a}\) denoting, as in § 1, no. 1, the set of \(x\in \mathbf{K}\) such that \(xa\in \mathbf{A}\) ).
\end{enumerate}

We know (Chapter II, § 2, no. 2, Corollary 2 to Proposition 4) that the condition \(\mathbf{M}_{\mathfrak{p}} = 0\) is equivalent to \(\mathfrak{a} \nsubseteq \mathfrak{p}\) and hence to \(\mathfrak{a}\mathbf{A}_{\mathfrak{p}} = \mathbf{A}_{\mathfrak{p}}\) (Chapter II, § 2, no. 5, Remark); on the other hand, for every integral ideal \(\mathfrak{b} \neq 0\) of A, the relation `` \(\mathfrak{b}\mathbf{A}_{\mathfrak{p}} = \mathbf{A}_{\mathfrak{p}}\) for all \(\mathfrak{p} \in \mathbb{P}\) '' is equivalent to \(\operatorname{div} \mathfrak{b} = \operatorname{div} \mathbf{A} = 0\) in D(A) (§ 1, no. 4, Proposition 7), or also to \(\operatorname{div}(\mathbf{A}: \mathfrak{b}) = 0\) and, as A: \(\mathfrak{b}\) is divisorial (§ 1, no. 1, Proposition 1), this relation is also equivalent to A: \(\mathfrak{b} = \mathbf{A}\) . The proposition then follows by noting that to say that \(\mathfrak{a} \nsubseteq \mathfrak{p}\) for \(\mathfrak{p} = (0)\) means that \(\mathfrak{a} \neq (0)\) .

Remark (1) The equivalent conditions of Proposition 9 mean also that \(\operatorname{Ass}(\mathbf{M})\) contains no prime ideal of height \(\leqslant 1\) . * They may also be interpreted by saying that \(\operatorname{Supp}(\mathbf{M})\) is of codimension \(\geqslant 2\) in \(\operatorname{Spec}(\mathbf{A})\) . *

DEFINITION 2. An A-module M is called pseudo-zero if it is finitely generated and it satisfies the equivalent conditions of Proposition 9.

This definition and Proposition 9 show that a pseudo-zero A-module is a torsion A-module; the converse is false.

\textbf{Examples}\par

\begin{enumerate}
\def\labelenumi{(\arabic{enumi})}
\item
  If A is a Dedekind domain, every prime ideal of A is of height \(\leqslant1\) ; to say that M is pseudo-zero means then that Supp(M) = \(\varnothing\) and hence that M = 0 (Chapter II, § 4, no. 4).
\item
  Let \(k\) be a field and \(A = k[X, Y]\) the polynomial ring over \(k\) in two indeterminates; if \(m\) is the maximal ideal \(AX + AY\) of \(A\) , the \(A\) -module \(A/m\) is pseudo-zero; for its annihilator \(m\) is not of height \(\leqslant 1\) since it contains the principal prime ideals \(AX\) and \(AY\) and is distinct from them; therefore \(A: m = A\) (§ 1, no. 6, Corollary 1 to Theorem 3).
\end{enumerate}

DEFINITION 3. Let M and N be two A-modules and f: M → N a homomorphism. f is called pseudo-injective (resp. pseudo-surjective, pseudo-zero) if Ker(f) (resp. Coker(f), Im(f)) is pseudo-zero; f is called pseudo-bijective if it is both pseudo-injective and pseudo-surjective.

A pseudo-bijective homomorphism is also called a pseudo-isomorphism.

Suppose that M and N are finitely generated; then, for \(f: M \to N\) to be pseudo-injective (resp. pseudo-surjective, pseudo-zero), it is necessary and sufficient that, for all \(p \in P \cup \{\{0\}\}, f_{p}: M_{p} \to N_{p}\) be injective (resp. surjective, zero); this follows from the flatness of the A-module \(A_{p}\) (cf.~Chapter I, § 2, no. 3, Remark 2).

Example (3) Let M be a torsion-free finitely generated A-module; then the canonical mapping \(c_{M}: M \to M^{**}\) of M to its bidual is a pseudo-isomorphism. For M is identified with a lattice of \(V = M \otimes_{A} K\) (no. 1, Proposition 1); we have seen that \(M_{p} = M_{p}^{**}\) for all \(p \in P\) (no. 2, Example 2) and, for p = 0, \(M_{p}\) and \(M_{p}^{**}\) are both equal to V.

THEOREM 4. Let E be a finitely generated A-module, T the torsion submodule of E and M = E/T. There exists a pseudo-isomorphism

\[
f \colon \mathbf {E} \rightarrow \mathbf {T} \times \mathbf {M}.
\]

We shall first show two lemmas.

LEMMA 2. Let \((\mathfrak{p}_i)_{1 \leqslant i \leqslant k}\) be a non-empty finite family of prime ideals of A of height 1 and let \(S = \bigcap_{i} (A - \mathfrak{p}_i)\) ; then the ring \(S^{-1}A\) is a principal ideal domain.

\(S^{-1}A\) is a semi-local ring whose maximal ideals are the \(m_{i} = p_{i}S^{-1}A\) for \(1 \leqslant i \leqslant k\) , the local ring \((S^{-1}A)_{m_{i}}\) being isomorphic to \(A_{p_{i}}\) (Chapter II, §3, no. 5, Proposition 17) and hence a discrete valuation ring. The ring \(S^{-1}A\) is therefore a Dedekind domain (§2, no. 2, Theorem 1 (f) and, as it is semi-local it is a principal ideal domain (§2, no. 2, Proposition 1).

LEMMA 3. There exists a homomorphism \(g: \mathbf{E} \to \mathbf{T}\) whose restriction to \(\mathbf{T}\) is both a homothety and a pseudo-isomorphism.

Let \(\mathfrak{a}\) be the annihilator of \(\mathrm{T}\) ; as \(\mathrm{T}\) is a finitely generated torsion A-module, \(\mathfrak{a} \neq 0\) . Let \(\mathfrak{p}_i\) ( \(1 \leqslant i \leqslant k\) ) be the prime ideals of height 1 containing \(\mathfrak{a}\) (which are finite in number (§ 1, no. 6, Theorem 4)); if this number is 0, T is pseudo-zero (Proposition 9(a)) and we may take g = 0. Otherwise, let \(S = \bigcap_{i} (A - p_{i})\) ; by Lemma 2, \(S^{-1}A\) is a principal ideal domain and hence \(S^{-1}M\) , which is a torsion-free finitely generated \(S^{-1}A\) -module, is free (Algebra, Chapter VII, § 4, no. 3, Corollary 2 to Theorem 2) and, as \(S^{-1}M = (S^{-1}E)/(S^{-1}T)\) , \(S^{-1}T\) is a direct factor of \(S^{-1}E\) (Algebra, Chapter II, § 1, no. 11, Proposition 21). Now,

\[
\operatorname{Hom} _ {\mathbf {S} ^ {- 1} \mathbf {A}} (\mathrm{S} ^ {- 1} \mathrm{E}, \mathrm{S} ^ {- 1} \mathrm{T}) = \mathrm{S} ^ {- 1} \operatorname{Hom} _ {\mathbf {A}} (\mathrm{E}, \mathrm{T})
\]

(Chapter II, §2, no. 7, Proposition 19); hence there exist \(s_0 \in S\) and \(g_0 \in \mathrm{Hom}_{\mathbf{A}}(\mathbf{E}, \mathbf{T})\) such that \(s_0^{-1}g_0\) is a projector of \(\mathbf{S}^{-1}\mathbf{E}\) onto \(\mathbf{S}^{-1}\mathbf{T}\) . If \(h_0 \in \mathrm{Hom}_{\mathbf{A}}(\mathbf{T}, \mathbf{T})\) denotes the restriction of \(g_0\) to \(\mathbf{T}\) , there therefore exists \(s_1 \in \mathbf{S}\) such that \(s_1h_0(x) = s_1s_0x\) for all \(x \in \mathbf{T}\) ; writing \(s = s_1s_0\) , \(g = s_1g_0\) , \(h = s_1h_0\) , \(h\) is therefore the homothety of ratio \(s\) on \(\mathbf{T}\) and is the restriction of \(g\) to \(\mathbf{T}\) . It remains to verify that \(h\) is a pseudo-isomorphism. Now, if \(p = 0\) or if \(p \in P\) is distinct from the \(p_i (1 \leqslant i \leqslant k)\) , \(T_p = 0\) (Chapter II, §4, no. 4, Proposition 17) and \(h_p: T_p \to T_p\) is an isomorphism; if on the contrary \(p\) is equal to one of the \(p_i (1 \leqslant i \leqslant k)\) , \(s\) is invertible in \(A_{p_i}\) and \(h_{p_i}\) , the homothety of ratio \(s\) on \(T_{p_i}\) , is also an isomorphism, which completes the proof of Lemma 3.

We now prove Theorem 4. Let \(g: \mathbf{E} \to \mathbf{T}\) be a homomorphism satisfying the properties of Lemma 3; let \(h\) be the restriction of \(g\) to \(\mathbf{T}\) and let \(\pi\) be the canonical projection of \(\mathbf{E}\) onto \(\mathbf{M}\) . We show that the homomorphism \(f = (g, \pi): \mathbf{E} \to \mathbf{T} \times \mathbf{M}\) solves the problem. There is the commutative diagram:

\[
\begin{array}{c c c} 0 \longrightarrow \mathrm{T} \longrightarrow & \mathrm{E} & \longrightarrow \mathrm{M} \longrightarrow 0 \\ h \Big \downarrow & f \Big \downarrow & ^ {1 _ {\mathrm{M}}} \Big \downarrow \\ 0 \longrightarrow \mathrm{T} \longrightarrow \mathrm{T} \times \mathrm{M} \longrightarrow \mathrm{M} \longrightarrow 0 \end{array}
\]

where the rows are exact. The snake diagram (Chapter I, § 1, no. 4, Proposition 2) gives the exact sequence:

\[
0 \rightarrow \operatorname{Ker} (h) \rightarrow \operatorname{Ker} (f) \rightarrow 0 \rightarrow \operatorname{Coker} (h) \rightarrow \operatorname{Coker} (f) \rightarrow 0
\]

and hence \(\operatorname{Ker}(f)\) is isomorphic to \(\operatorname{Ker}(h)\) and \(\operatorname{Coker}(f)\) to \(\operatorname{Coker}(h)\) . As \(h\) is a pseudo-isomorphism, so is \(f\) .

We can say that ``to within a pseudo-isomorphism'' Theorem 4 reduces the study of finitely generated A-modules to that of torsion-free modules on the one hand and to that of torsion modules on the other. Moreover, we have seen above (Example 3) that a torsion-free module is pseudo-isomorphic to its bidual and hence to a reflexive module. As for torsion modules, there is the following result, which determines them to within a pseudo-isomorphism:

THEOREM 5. Let T be a finitely generated torsion A-module. There exist two finite families \((n_i)_{i\in I}\) and \((p_i)_{i\in I}\) , where the \(n_i\) are integers \(\geqslant 1\) and the \(p_i\) are prime ideals of A of height 1 such that, if we write \(\mathrm{T}' = \bigoplus_{i\in I} \mathrm{A} / p_i^{n_i}\) , there exists a pseudo-isomorphism of T to \(\mathrm{T}'\) . Moreover, the families \((n_i)_{i\in I}\) and \((p_i)_{i\in I}\) with this property are unique to within a bijection of the indexing set and the \(p_i\) contain the annihilator of T.

Uniqueness: If \(f\colon T\to T'\) is a pseudo-isomorphism and \(\mathfrak{p}\in\mathbf{P},f_{\mathfrak{p}}\colon T_{\mathfrak{p}}\to T'_{\mathfrak{p}}\) is an isomorphism. Now, \(T'_{\mathfrak{p}}\) is the direct sum of the \(A_{\mathfrak{p}}/p^{n_{i}}A_{\mathfrak{p}}\) , the sum being over the indices \(i\) such that \(\mathfrak{p}_{i}=\mathfrak{p}\) ; the \(p^{n_{i}}A_{\mathfrak{p}}\) are therefore the elementary divisors of the torsion \(A_{\mathfrak{p}}\) -module \(T_{\mathfrak{p}}\) (Algebra, Chapter VII, § 4, no. 7); their uniqueness has been proved in Algebra, Chapter VII, § 4, no. 7, Proposition 7.

Existence: We may confine our attention to the case where \(\mathbf{T} \neq 0\) . Let \(\mathfrak{a}\) be the annihilator (non-zero and distinct from A) of \(\mathbf{T}\) , \(\mathfrak{p}_i\) ( \(1 \leqslant i \leqslant k\) ) the prime ideals of A of height 1 containing \(\mathfrak{a}\) (which are finite in number (§ 1, no. 6, Theorem 4)) and \(\mathbf{S} = \bigcap_i (\mathbf{A} - \mathfrak{p}_i)\) . The semi-local ring \(\mathbf{A}' = \mathbf{S}^{-1}\mathbf{A}\) is a principal ideal domain (Lemma 2) and has maximal ideals the \(\mathfrak{m}_i = \mathfrak{p}_i\mathbf{A}'\) ; as \(\mathbf{S}^{-1}\mathbf{T}\) is a finitely generated torsion \(\mathbf{A}'\) -module, it is isomorphic to a finite direct sum \(\bigoplus_{f \in \mathbf{I}} \mathbf{A}' / \mathfrak{m}_{\phi(f)}^{n_f}\) , where \(\phi\) is a mapping of a finite set I to \([1, k]\) (Algebra, Chapter VII, § 4, no. 7, Proposition 7); as \(\mathbf{A}' / \mathfrak{m}_{\phi(f)}^{n_f}\) is isomorphic to \(\mathbf{S}^{-1}(\mathbf{A} / \mathfrak{p}_{\phi(f)}^{n_f})\) (Chapter II, § 2, no. 4), we have obtained a torsion A-module \(\mathbf{T}'\) of the desired type and an isomorphism \(f_0\) of \(\mathbf{S}^{-1}\mathbf{T}\) onto \(\mathbf{S}^{-1}\mathbf{T}'\) . As \(\operatorname{Hom}_{\mathbf{S}^{-1}\mathbf{A}}(\mathbf{S}^{-1}\mathbf{T}, \mathbf{S}^{-1}\mathbf{T}')\) is equal to \(\mathbf{S}^{-1}\operatorname{Hom}_{\mathbf{A}}(\mathbf{T}, \mathbf{T}')\) (Chapter II, § 2, no. 7, Proposition 19), there exist \(s \in S\) and a homomorphism \(f: T \to T'\) such that \(f_0 = s^{-1}f\) . It remains to show that \(f\) is a pseudo-isomorphism: now, if \(\mathfrak{p} = 0\) or if \(\mathfrak{p} \in P\) is distinct from the \(\mathfrak{p}_i\) , \(\mathrm{T}_{\mathfrak{p}} = \mathrm{T}_{\mathfrak{p}}' = 0\) (Chapter II, § 4, no. 4, Proposition 17); if on the contrary \(\mathfrak{p}\) is one of the \(\mathfrak{p}_i\) ( \(1 \leqslant i \leqslant k\) ), \(s\) is invertible in \(\mathbf{A}_{\mathfrak{p}_i}\) and, as \(f_{\mathfrak{p}_i} = s(f_0)_{\mathfrak{p}_i}\) and \((f_0)_{\mathfrak{p}_i}\) is an isomorphism of \(\mathrm{T}_{\mathfrak{p}_i} = (\mathrm{S}^{-1}\mathrm{T})_{\mathfrak{m}_i}\) onto \(\mathrm{T}_{\mathfrak{p}_i}' = (\mathrm{S}^{-1}\mathrm{T}')_{\mathfrak{m}_i}\) , so is \(f_{\mathfrak{p}_i}\) .

Remark (2) In the statement of Theorem 5, the modules \(\mathbf{A} / \mathfrak{p}_i^n\) may be replaced by \(\mathbf{A} / \mathfrak{p}_i^{(n_i)}\) (§ 1, no. 4, Proposition 8). For all \(\mathfrak{p} \in \mathbb{P}\) , the canonical mapping \(g\colon \mathbf{A} / \mathfrak{p}^n \to \mathbf{A} / \mathfrak{p}^{(n)} = \mathbf{A} / (\mathbf{A} \cap \mathfrak{p}^n \mathbf{A}_{\mathfrak{p}})\) is a pseudo-isomorphism, as, for all \(\mathfrak{q} \in \mathbb{P}\) distinct from \(\mathfrak{p}\) , \(\mathbf{A}_{\mathfrak{q}} / \mathfrak{p}^n \mathbf{A}_{\mathfrak{q}} = \mathbf{A}_{\mathfrak{q}} / \mathfrak{p}^{(n)} \mathbf{A}_{\mathfrak{q}} = 0\) and \(\mathbf{A}_{\mathfrak{p}} / \mathfrak{p}^n \mathbf{A}_{\mathfrak{p}} = \mathbf{A}_{\mathfrak{p}} / \mathfrak{p}^{(n)} \mathbf{A}_{\mathfrak{p}}\) .

* Given an exact sequence of A-modules, E → F → G, if E and G are pseudo-zero, so is F, as follows from Definition 2 and Chapter II, § 2, no. 4, Theorem 1. In the language of categories, we may then say that, in the category C of A-modules, the sub-category C' of pseudo-zero modules is full and we may then define the quotient category C/C': the objects in this category are also A-modules but the set of morphisms from E to F (for E, F in C/C') is the direct limit of the set of commutative groups \(\mathrm{Hom}_{\mathbf{A}}(\mathbf{E}',\mathbf{F}')\) , where \(\mathbf{E}'\) (resp. \(\mathbf{F}'\) ) runs through the set of submodules of \(\mathbf{E}\) (resp. the set of quotient modules \(\mathbf{F} / \mathbf{F}''\) of \(\mathbf{F}\) ) such that \(\mathbf{E} / \mathbf{E}'\) (resp. \(\mathbf{F}''\) ) is pseudo-zero. Of course, for every ordered pair of A-modules \(\mathbf{E}\) , \(\mathbf{F}\) , there is a canonical homomorphism \(\mathrm{Hom}_{\mathcal{C}}(\mathbf{E},\mathbf{F})\to \mathrm{Hom}_{\mathcal{C} / \mathcal{C}'}(\mathbf{E},\mathbf{F})\) . To say that a homomorphism \(u\in \mathrm{Hom}_{\mathbf{A}}(\mathbf{E},\mathbf{F})\) is pseudo-zero (resp. pseudo-injective, pseudo-surjective, pseudo-bijective) means that its canonical image in \(\mathrm{Hom}_{\mathcal{C} / \mathcal{C}'}(\mathbf{E},\mathbf{F})\) is zero (resp. a monomorphism, an epimorphism, an isomorphism).*

